\section{Introduzione}

\subsection{Il problema che non scompare}
Perché sente il dolore? È la domanda più semplice del mondo. E al tempo stesso la più difficile. Sappiamo tutti che cos'è il dolore. Lo abbiamo sentito tutti. Ma perché lo sentiamo? Sappiamo che i nocicettori inviano segnali al talamo, che la corteccia li elabora e che vengono rilasciate endorfine. Tutto questo è una descrizione funzionale. Ci dice che cosa fa il cervello, ma non ci dice perché lei lo sente.

Pensi a un computer. Può elaborare informazioni, riconoscere un'immagine, giocare a scacchi. Ma non sente nulla. Non ha esperienza soggettiva. Nulla gli fa male. Perché il cervello è diverso? Questo è il problema difficile della coscienza, formulato da David Chalmers nel 1995. Da allora le neuroscienze hanno fatto enormi progressi, ma il problema difficile è rimasto intatto. Nessuno è riuscito a spiegare perché la materia diventa esperienza. Nel 1998 Christof Koch scommise con Chalmers che un correlato neurale della coscienza sarebbe stato identificato entro 25 anni. Nel 2023 ha ammesso che la scommessa era persa.

\subsection{La materia oscura --- le cifre esatte}
L'universo ha una composizione che conosciamo abbastanza bene. La missione Planck (2018) ci ha dato le cifre esatte. Solo il 4,9\% dell'universo è materia ordinaria --- atomi, stelle, pianeti, persone. La materia oscura rappresenta il 26,8\% e l'energia oscura il 68,3\%. In pratica, il 95,1\% dell'universo è sconosciuto.

Immagini l'universo come un forziere con 100 monete. Solo 5 monete sono visibili. Le altre 95 sono invisibili. Di queste, 27 sono materia oscura e 68 energia oscura. Della materia oscura sappiamo che esiste (le prove gravitazionali sono solide), che non emette né assorbe luce, e che ci attraversa senza interagire. Circa 100 milioni di particelle attraversano il suo corpo ogni secondo. Non sappiamo di che cosa sia fatta.

Una restrizione recente (2025), dovuta a Bonvin et al., ha dimostrato che la materia oscura cade nei pozzi gravitazionali esattamente come la materia ordinaria. Se esiste una «quinta forza», questa non può superare il 7\% della forza gravitazionale.

\subsection{Che cos'è ART-12D?}
ART-12D parte da un'idea semplice: l'universo ha più dimensioni delle quattro che percepiamo. Queste dimensioni aggiuntive sono compatte --- così piccole che non le notiamo. La materia oscura esiste in tutte e dodici le dimensioni. Proiettata in 4D, appare come un campo che riempie l'universo. Questo campo è ciò che il cervello riceve come coscienza.

Il nome deriva da: \textbf{Akashic} (il campo informazionale universale), \textbf{Resonance} (perché l'interazione è una risonanza), \textbf{Theory} (un quadro teorico) e \textbf{12D} (perché lavoro in dodici dimensioni).

I tre postulati sono: 1. La materia oscura è un campo quantistico fondamentale. 2. Il cervello è un ricevitore geometrico. 3. L'accoppiamento avviene tramite interazione di Yukawa.

Pensi a una radio. Non crea la musica, la capta dall'aria. Le onde radio esistono che la radio sia accesa o no. Quando la radio si guasta, la musica non scompare dall'universo --- cessa solo la sua ricezione locale. Allo stesso modo, il cervello funziona esattamente come una radio. La coscienza non è generata dai neuroni, è ricevuta da essi. E ciò che ricevono è il campo della materia oscura.

\subsection{Perché dodici dimensioni?}
Perché dodici? Perché non 10, o 11, o 26? La risposta viene dalla teoria delle stringhe. Immagini una corda di chitarra. Può vibrare in più modi. Ogni modo corrisponde a una nota. Allo stesso modo, una stringa fondamentale può vibrare in più dimensioni. Il numero di dimensioni determina quali tipi di vibrazione sono possibili.

La storia delle dimensioni: Kaluza--Klein (5 dimensioni, 1919--1926), superstringhe (10, 1970--1980), teoria M (11, Witten, 1995), teoria F (12, Vafa, 1996), ART-12D (12, Ofiteru, 2026). Dodici è il maggior numero di dimensioni in cui la teoria delle stringhe è matematicamente coerente. Più dimensioni porterebbero a incoerenze; meno a una teoria incompleta.

\subsection{Come sono le dimensioni compatte?}
Immagini un tubo da giardino. Da lontano sembra una linea 1D --- ha solo lunghezza. Da vicino ha spessore --- è 2D (la superficie del cilindro). Allo stesso modo, il nostro spazio sembra tridimensionale. Ma da vicino (alla scala di Planck, $10^{-35}$ m) potrebbe avere undici dimensioni. Le otto dimensioni compatte sono arrotolate in ogni punto dello spazio.

Che cosa significa «arrotolate»? Non sono luoghi dove si va. Non sono mondi paralleli che si raggiungono viaggiando. Sono proprietà locali di ogni punto. Ogni punto dello spazio-tempo ha, in più, otto piccole dimensioni «attaccate» ad esso.

\subsection{Come si legano le dimensioni alla coscienza?}
La materia oscura esiste in dodici dimensioni. È un campo scalare $A(x,y)$, dove $x$ sono le coordinate 4D e $y$ le otto coordinate compatte. Quando compattifichiamo le otto dimensioni, il campo si decompone in modi di Kaluza--Klein:
\begin{equation}
A(x,y) = \sum_n A_n(x) Y_n(y)
\end{equation}
Ogni modo ha una massa e una frequenza caratteristica:
\begin{equation}
m_n^2 = \frac{n^2}{R^2}
\end{equation}
Il cervello, attraverso la sua struttura geometrica (i microtubuli nei neuroni), può risuonare con certi modi. Quando risuona, il campo viene trasdotto in esperienza soggettiva.

Perché conta l'architettura e non la dimensione? Il cervelletto ha 69 miliardi di neuroni (l'80\% del totale), ma non produce coscienza. Ha un'architettura \textit{feedforward} --- l'informazione scorre in una sola direzione, come un'antenna a filo. La corteccia ha solo 16 miliardi di neuroni, ma è la sede della coscienza. Ha un'architettura \textit{ricorrente} --- l'informazione circola in anelli, come un'antenna parabolica ben orientata. Conclusione: non conta il numero di neuroni, ma la geometria della rete.

\section{Metodologia}
\subsection{Come ho lavorato}
Non sono un fisico professionista. Non ho un laboratorio. Non ho finanziamenti. Quello che ho è una biblioteca, tempo e una domanda che non mi lascia in pace. Ho lavorato in quattro passi. Il primo passo è stato la formulazione dell'ipotesi, partendo dall'intuizione che la coscienza è ricevuta e non generata. Il secondo passo è stato la formalizzazione, con gli strumenti della matematica nota: Kaluza--Klein per le dimensioni aggiuntive, la teoria F per dodici dimensioni, la teoria dei campi per le lagrangiane, l'IIT per l'informazione integrata. Il terzo passo è stato la derivazione delle predizioni, cioè ciò che dovremmo osservare se la teoria è corretta. Ho selezionato le predizioni testabili oggi. Il quarto passo è stato la proposta di esperimenti con mezzi modesti, a cominciare da uno da 50.000 dollari.

\subsection{Che cosa mi ha guidato}
Mi sono attenuto a quattro criteri. La coerenza interna --- le equazioni non devono contraddirsi. La falsificabilità --- la teoria deve fare predizioni testabili. La convergenza --- la teoria non deve contraddire la fisica consolidata. L'onestà --- devo riconoscere che è speculativa.

\subsection{Che cosa ho usato}
Ho usato Kaluza--Klein (1919) per l'idea di dimensioni compatte, la teoria F (Vafa, 1996) per il quadro a dodici dimensioni, l'IIT (Tononi, 2004) per la misura dell'informazione integrata, Orch OR (Penrose, 1996) per il collasso quantistico nei microtubuli, e la teoria quantistica dei campi per le lagrangiane e le equazioni del moto.

\subsection{I miei limiti}
Devo essere onesto sui limiti. Non ho prove sperimentali; la teoria non è stata testata. Non ho revisione paritaria; non ho pubblicato su riviste scientifiche. Non ho affiliazione accademica; sono un ricercatore indipendente. Non ho finanziamenti; i miei mezzi sono limitati.
